\documentclass[a4paper]{article}
\usepackage{amsfonts}
\usepackage{amsmath}
\usepackage{amssymb}
\usepackage{graphicx}	
\newcommand{\integraal}{\displaystyle\int}

\begin{document}
  
\noindent \large Auteur: Joe Ayoub \\
\noindent \large Studierichting: Informatica\\
\noindent \large Oefeningenreeks 4A nr.42.18 \\

\medskip

\normalsize

$\integraal{}\dfrac{2x^2 - 5x + 7}{x^3 - 7x + 6}dx = \integraal{}\dfrac{2x^2 - 5x + 7}{(x-1)(x-2)(x+3)}dx$ \\

$ \dfrac{2x^2 - 5x + 7}{(x-1)(x-2)(x+3)}dx = \dfrac{A}{(x-1)} + \dfrac{B}{(x-2)} + \dfrac{C}{(x + 3)}$ \\

Door gebruik te maken van de Fuss methode kunnen we nu zeggen dat: \\

$ \Rightarrow A = \dfrac{2x^2 - 5x + 7}{(x-2)(x+3)}$ \\

Als $x = $ dan $A =  -1$ \\

$ \Rightarrow B = \dfrac{2x^2 - 5x + 7}{(x-1)(x+3)}$\\

Als $x = 2$ dan $B = 1$ \\

$ \Rightarrow C = \dfrac{2x^2 - 5x + 7}{(x-1)(x-2)}$\\

Als $x = -3$ dan $C = 2$ \\

We gaan terug naar het integraal \\

$= \integraal{}\dfrac{2x^2 - 5x + 7}{(x-1)(x-2)(x+3)}dx$ \\

$= -\integraal{}\dfrac{1}{x-1}dx + \integraal{}\dfrac{1}{x-2}dx + 2\integraal{}\dfrac{1}{x+3}dx$ \\

$= -\ln{}|x-1| + \ln{}|x-2| + 2\ln{}|x+3| + C$

\begin{thebibliography}{999}
%\bibitem{a} Referentie
\end{thebibliography}
\end{document} 
